\chapter{知识的融合}
\label{chap:knowledge-fusion}

前几章得到的知识可能来自人工、规则、模型或数据本身。同一个样本因此可能同时收到多份答案：三名标注者给出两个类别，两条规则复用了同一词典，一个模型只对容易样本作答。答案多并不等于证据充分；若它们描述的不是同一对象，或重复使用同一信息，直接计票反而会放大错误。本章要回答的问题是：怎样把这些来源不同、质量不同且可能冲突的知识组织成可解释的综合结果？

融合之前需要解决三个层次的问题。来源建模描述每个来源会在何处作答、容易怎样出错以及是否共享信息；知识对齐确定各条记录是不是在讨论同一对象和同一命题；融合推断再把可比较的证据组合起来。推断结果也不必总是一个硬标签：证据不足时可以保留概率分布、候选集合或弃权，语义不同的知识则应保留各自目标并联合使用。下面沿用一组人工、规则和模型产生的二分类标签，逐步说明这条处理链。

\section{知识来源的建模}
\label{sec:knowledge-fusion-source-modeling}

\term{知识来源}（knowledge source）是能够为一批对象产生知识记录的主体或过程，例如一名标注者、一条规则、一个模型，或同一模型的一种固定配置。划分来源时不能只看输出列名：两条规则若来自同一模板或词典，两个模型若共享训练数据和基础模型，这些共享条件也属于来源记录。来源模型的输入包括来源身份、覆盖状态、已对齐的输出和已知共享关系，输出则是后续融合需要的能力、混淆与依赖参数。

来源的\term{可靠性}（reliability）概括“这个来源的判断有多可信”，\term{混淆}（confusion）进一步回答“真实类别为某一类时，它会错向哪一类”，\term{相关性}（correlation）则描述“多个来源是否重复使用了同一信息”。可靠性是边缘能力，混淆保留错误方向，相关性考察来源之间的联合行为。三者不能互相替代：两个总体正确率相同的来源可能有不同混淆，四个高正确率来源也可能只是同一个错误机制的四份副本。

\subsection{来源可靠性的建模}
\label{sec:knowledge-fusion-reliability}

设有$n$个已经完成对象与语义对齐的样本、$m$个来源和类别集合$Y=\{1,\ldots,K\}$。来源$s$对样本$i$的有效输出记为$\lambda_{i,s}\in Y$，$o_{i,s}=1$表示给出了有效类别，$o_{i,s}=0$表示没有有效答案。没有答案还要继续区分未请求、执行失败和主动弃权；指标中的$o_{i,s}$只负责说明当前是否存在可比较标签，不抹去缺失原因。

\subsubsection{利用参考标签估计可靠性}
\label{sec:knowledge-fusion-reference-reliability}

若一部分样本有经过独立审计的参考标签，就可以直接估计来源能力。记审计样本索引集合为$I_{\mathrm{audit}}$，参考标签为$y_i$。来源$s$的覆盖率与覆盖条件下正确率分别为
\begin{equation}
  \widehat c_s
  =\frac{1}{n}\sum_{i=1}^{n}o_{i,s},
  \qquad
  \widehat a_s
  =\frac{\sum_{i\in I_{\mathrm{audit}}}
    o_{i,s}\mathbb{I}\{\lambda_{i,s}=y_i\}}
  {\sum_{i\in I_{\mathrm{audit}}}o_{i,s}}.
  \label{eq:knowledge-fusion-audited-reliability}
\end{equation}
第一个分母包含全部目标样本，回答来源在多大范围内给出答案；第二个分母只包含审计集合中该来源实际作答的样本，回答它作答后有多大比例与参考一致。将未请求或弃权直接计作错误会混合覆盖与正确性，将弃权计作负类则会改变任务语义。

表~\ref{tab:knowledge-fusion-source-records}给出本章使用的人工教学记录。规则$R_1$与$R_2$来自同一模板，人工来源$H$独立作答，模型$M$可以弃权；前四个样本有独立参考标签，后四个只用于展示融合。所有数值均为教学设定，不是实验结果。

\begin{table*}[htbp]
  \centering
  \small
  \begin{tabularx}{\linewidth}{@{}clccccc@{}}
    \toprule
    \textbf{样本} & \textbf{区域} & \textbf{参考} & \textbf{$R_1$} & \textbf{$R_2$} & \textbf{$H$} & \textbf{$M$} \\
    \midrule
    $x_1$ & 简单 & 甲 & 甲 & 甲 & 甲 & 甲 \\
    $x_2$ & 困难 & 乙 & 甲 & 甲 & $\varnothing_{\text{未请求}}$ & 乙 \\
    $x_3$ & 简单 & 甲 & 甲 & 甲 & 甲 & 乙 \\
    $x_4$ & 困难 & 乙 & $\varnothing_{\text{失败}}$ & $\varnothing_{\text{失败}}$ & $\varnothing_{\text{未请求}}$ & 甲 \\
    \midrule
    $x_5$ & 简单 & 未审计 & 甲 & 甲 & 乙 & $\bot$ \\
    $x_6$ & 简单 & 未审计 & 甲 & 甲 & 乙 & 乙 \\
    $x_7$ & 未覆盖 & 未审计 & $\varnothing_{\text{失败}}$ & $\varnothing_{\text{失败}}$ & $\varnothing_{\text{未请求}}$ & $\bot$ \\
    $x_8$ & 简单 & 未审计 & 乙 & 乙 & 甲 & 甲 \\
    \bottomrule
  \end{tabularx}
  \caption{多来源标签与覆盖记录。$\varnothing_{\text{未请求}}$表示没有向来源请求答案，$\varnothing_{\text{失败}}$表示请求后未完成，$\bot$表示来源运行后主动弃权；前三种状态都不作为类别票。$x_1$至$x_4$参与参考审计，人工来源$H$只处理简单区，模型$M$同时覆盖简单区与困难区；所有对象和类别语义均假设已经对齐。}
  \label{tab:knowledge-fusion-source-records}
\end{table*}
\FloatBarrier

$R_1$、$R_2$与$M$在整张表上的有效覆盖率都是$6/8$，只处理简单区的$H$为$5/8$。在审计记录中，$R_1$与$R_2$各答对$2/3$，$H$答对$2/2$，覆盖两个区域的$M$答对$2/4$；若只比较这些原始比例，会得到$H$优于$M$的印象。限制到二者都覆盖的简单审计区$x_1,x_3$后，$H$为$2/2$，$M$为$1/2$；困难审计区$x_2,x_4$则只有$M$的$1/2$，$H$根本没有可估计结果。因而$H$在已覆盖简单样本上的高正确率不能推出它在困难区域同样可靠。对区域$g$的正确率应把式~\eqref{eq:knowledge-fusion-audited-reliability}的分子和分母同时限制到$I_g$，并报告分母；分组越细，需要的审计样本也越多。

参考标签并非天然无误。审计样本的抽取方式、参考答案的形成过程和适用版本都应与来源的目标使用范围对应。参考数量太少时应保留区间或后验不确定性，某一区域完全没有参考时应标为未估计，而不是沿用总体平均。来源会选择性作答时，还要单独检查$o_{i,s}$与类别、难度和时间的关系，否则观察到的高正确率可能只是“只答容易题”的结果。

\subsubsection{利用来源关系估计可靠性}
\label{sec:knowledge-fusion-relational-reliability}

没有参考标签时，只能从来源在重叠样本上的赞同与分歧间接估计能力。这里真正可观察的是标签矩阵$\bm{\Lambda}=(\lambda_{i,s})$及其覆盖掩码$\bm{O}$，待估计的则是来源参数、类别比例和每个样本的潜在答案。多来源概率模型把这些未知量联系起来：若干相互独立且方向已知的来源反复赞同时，一个共同潜在类别可以解释这种一致；若来源经常在某一类别上系统分歧，类别相关误差又能给出另一种解释。

这类估计的关键不是优化器能否输出一个数，而是模型是否\term{可识别}（identifiable）：不同参数解释能否产生完全相同的观测分布。若两个来源从不覆盖同一批样本，它们的相对能力缺少比较桥梁；若十个来源复制同一规则，它们的十列输出只提供一类信息；若交换两个潜在类别的名称并同时交换类别先验和混淆矩阵的行，观测标签的分布仍然不变。三种困难分别来自重叠不足、有效独立来源不足和类别置换，不能归结为同一种“样本太少”。

设参数$\bm{\theta}$诱导观测分布$p_{\bm{\theta}}(\bm{\Lambda})$。严格可识别要求
\[
  p_{\bm{\theta}}(\bm{\Lambda})
  =p_{\bm{\theta}'}(\bm{\Lambda})
  \quad\Longrightarrow\quad
  \bm{\theta}=\bm{\theta}',
\]
或只允许预先声明的等价变换，例如潜在类别名称置换。为写清一类常用充分条件，把来源划分为三个非空组$G_1,G_2,G_3$，并假设三组联合输出在给定$Y$后条件独立。令$\Omega_j$为第$j$组所有来源的联合输出空间，包含模型显式保留的弃权状态；定义条件概率矩阵$\bm{M}_j\in\mathbb{R}^{K\times|\Omega_j|}$为
\[
  [\bm{M}_j]_{k,u}
  =\mathbb{P}(\bm{\Lambda}_{G_j}=u\mid Y=k),
  \qquad
  k\in Y,\ u\in\Omega_j.
\]
矩阵第$k$行因此是潜在类别$k$下该组联合输出的条件分布。这里的行Kruskal秩$\operatorname{krank}(\bm{M}_j)$是“任取$r$行都线性无关”的最大$r$。若
\[
  \operatorname{krank}(\bm{M}_1)
  +\operatorname{krank}(\bm{M}_2)
  +\operatorname{krank}(\bm{M}_3)
  \geq 2K+2,
\]
则相应潜类参数在一般位置上除类别置换外可识别\scite{allman2009identifiability} 这不是“只要有三个来源就足够”的保证：组内来源完全相同、恒定输出或退化混淆都会降低联合输出矩阵的秩，组间条件依赖还会直接破坏上述分解前提。

类别方向仍需外部锚点确定。少量独立审计样本可以直接命名潜在类别；也可以加入“多数来源平均优于随机”等方向假设，但这增加的是模型约束，不是验证证据。减少参数、合并相似来源或设置先验能够缓解有限样本估计，却不能把原本不可区分的现实解释自动变成已核实真值。后文的Dawid--Skene模型和标签模型会把这些条件落实到具体似然中。

\subsection{来源混淆与偏差的建模}
\label{sec:knowledge-fusion-confusion}

单一正确率把所有错误合并成一个数。多类别任务中，来源可能很会识别甲类，却经常把乙类误成丙类；即使两个来源总体正确率相同，这种方向差异也会让它们对当前冲突提供不同证据。对来源$s$定义混淆矩阵
\begin{equation}
  \theta^{(s)}_{k,\ell}
  =\mathbb{P}(\lambda_{i,s}=\ell
    \mid Y_i=k,o_{i,s}=1),
  \qquad
  \sum_{\ell=1}^{K}\theta^{(s)}_{k,\ell}=1.
  \label{eq:knowledge-fusion-confusion-matrix}
\end{equation}
行索引$k$是真实类别，列索引$\ell$是来源输出。对角项表示分类别正确率，非对角项指出错误流向；这一行列约定必须在实现、图表与公式中保持一致。

有参考标签时，可以在每个真实类别内计数并按行归一化。若来源$s$在参考为$k$的样本上输出$\ell$共$N^{(s)}_{k,\ell}$次，则极大似然估计为$\widehat\theta^{(s)}_{k,\ell}=N^{(s)}_{k,\ell}/\sum_{r}N^{(s)}_{k,r}$。少数类别的分母很小时，矩阵会十分不稳定，可以通过Dirichlet平滑、合并过细类别或报告区间降低误读风险，但平滑结果仍受先验影响。

考虑一个三分类教学例子：两个来源的总体正确率都为$0.8$，来源A的错误主要把乙类标成甲类，来源B的错误却主要把乙类标成丙类。它们在总体可靠性上相同，对“真实类别是否为乙”的反向证据却指向不同候选；当A输出甲、B输出丙时，单一正确率无法解释这次分歧，两个混淆矩阵的乙类行才能给出相应似然。这里的$0.8$是人为设定，用于说明错误方向，不是经验结论。

覆盖选择、类别偏好和类别混淆描述不同对象。来源频繁弃权属于覆盖机制；在真实类别相同的条件下偏向输出某一列，才体现在混淆矩阵中；数据本身的类别比例变化则改变先验$\mathbb{P}(Y=k)$。部署期观察到输出“甲”的比例升高，可能来自样本中甲类变多，也可能来自来源更偏向甲，单凭输出边缘分布无法判定来源已经退化。

全局混淆矩阵假设来源行为在所有样本上稳定。若能力随领域、难度或时间改变，可以估计$\theta^{(s,g)}_{k,\ell}$，或令混淆概率依赖样本特征。模型越细，参数越多，每个参数得到的重叠或参考证据越少。因而应先由使用范围决定必要分组，再检查每一组的有效样本量和更新频率，不能用任意细分制造看似精确的能力画像。弃权可以作为额外观察状态进入生成模型，也可以由独立覆盖模型描述；无论采用哪一种，都不能把弃权设成新的真实类别。

\subsection{来源相关性的建模}
\label{sec:knowledge-fusion-correlation}

来源相关性关心多份输出是否携带重复信息。规则$R_1$与$R_2$来自同一模板，即使它们在表~\ref{tab:knowledge-fusion-source-records}中两次投甲，也不能自动视作两份独立证据。类似共享还可能来自同一训练集、基础模型、检索结果、提示模板或标注指南。来源目录应先记录这些已知关系，再用对齐后的输出检查模型是否还遗漏了依赖。

直接计算两个来源的一致率不能区分两种机制：它们可能都正确地反映真实类别，也可能在给定真实类别后仍共享同一种错误。后者才是融合时重复计证据的关键。对两个始终作答的来源$s,t$，条件独立假设写作
\begin{equation}
  \mathbb{P}(\lambda_{i,s},\lambda_{i,t}\mid Y_i)
  =
  \mathbb{P}(\lambda_{i,s}\mid Y_i)
  \mathbb{P}(\lambda_{i,t}\mid Y_i).
  \label{eq:knowledge-fusion-conditional-independence}
\end{equation}
若等式不成立，单独的混淆矩阵无法完整描述联合误差。没有真值或可靠潜在标签时，所谓“残余相关性”仍依赖当前来源模型，不能仅凭一个经验相关系数确定共享机制或因果方向。

图~\ref{fig:knowledge-fusion-source-dependence}对比两种机制。左侧假设给定真实类别后各来源独立；右侧增加共享规则或证据节点，使三条规则在真实类别之外还受到共同影响。同样得到多张赞同票时，右侧的有效信息量可能更接近“一组重复规则加一份独立判断”，而不是四份独立判断。

\begin{figure*}[htbp]
  \centering
  \begin{tikzpicture}[
  x=1mm,y=1mm,
  every node/.style={font=\figurefont\small,text=BookInk},
  truth/.style={circle,draw=BookTeal,fill=FigureModel!12,line width=1pt,minimum size=11mm,inner sep=1pt},
  source/.style={book node,draw=BookBlue,fill=FigureInput!9,line width=0.8pt,minimum width=18mm,minimum height=9mm,inner sep=2pt},
  shared/.style={book node,draw=BookAmber,fill=BookAmber!10,line width=0.8pt,minimum width=24mm,minimum height=9mm,inner sep=2pt}
]
  \node[font=\figurefont\bfseries] at (35,26) {给定类别后独立};
  \node[truth] (yl) at (35,10) {$Y$};
  \node[source] (l1) at (8,-10) {来源1};
  \node[source] (l2) at (35,-10) {来源2};
  \node[source] (l3) at (62,-10) {来源3};
  \draw[book arrow,knowledge arrow,draw=BookTeal,line width=1pt] (yl) -- (l1);
  \draw[book arrow,knowledge arrow,draw=BookTeal,line width=1pt] (yl) -- (l2);
  \draw[book arrow,knowledge arrow,draw=BookTeal,line width=1pt] (yl) -- (l3);
  \node[font=\figurefont\scriptsize,text=BookMuted,align=center] at (35,-26) {三条输出路径分别提供信息};

  \draw[BookRule,line width=0.7pt] (73,-40) -- (73,29);

  \node[font=\figurefont\bfseries] at (124,26) {共享机制引入相关错误};
  \node[truth] (yr) at (124,10) {$Y$};
  \node[source] (r1) at (86,-10) {规则1};
  \node[source] (r2) at (109,-10) {规则2};
  \node[source] (r3) at (132,-10) {规则3};
  \node[source] (human) at (155,-10) {标注者$H$};
  \node[shared] (u) at (109,-38) {共享模板或证据};
  \draw[book arrow,knowledge arrow,draw=BookTeal,line width=1pt] (yr) -- (r1);
  \draw[book arrow,knowledge arrow,draw=BookTeal,line width=1pt] (yr) -- (r2);
  \draw[book arrow,knowledge arrow,draw=BookTeal,line width=1pt] (yr) -- (r3);
  \draw[book arrow,knowledge arrow,draw=BookTeal,line width=1pt] (yr) -- (human);
  \draw[book arrow,knowledge arrow,draw=BookAmber,dashed] (u.north west) -- (r1.south);
  \draw[book arrow,knowledge arrow,draw=BookAmber,dashed] (u.north) -- (r2.south);
  \draw[book arrow,knowledge arrow,draw=BookAmber,dashed] (u.north east) -- (r3.south);
  \draw[BookAmber,dashed,line width=0.8pt,rounded corners=1.2mm] (75,-21) rectangle (143,0);
  \node[font=\figurefont\scriptsize,text=BookAmber,anchor=south west] at (78,3) {同一来源组};
  \node[font=\figurefont\scriptsize,text=BookAmber,anchor=west] at (78,-28) {共享影响};
\end{tikzpicture}

  \caption{独立信息与共享错误来源。左面板表示给定真实类别后来源独立的模型；右面板包含三条同模板规则与一名独立标注者，虚线和分组框标出三条规则的共享影响。图展示两种待检验的模型假设，不表示可以只从输出一致率唯一恢复真实依赖结构。}
  \label{fig:knowledge-fusion-source-dependence}
\end{figure*}
\FloatBarrier

沿用右面板，设三条规则都输出甲，独立标注者$H$输出乙。保留四票时结果是甲的$3:1$多数；把三条完全重复的规则合成一个来源组后，结果变成甲、乙各一票的未决。依赖模型则不必强行合票：例如在平衡先验的教学设定中，把规则组作为正确率$0.7$的一个联合因子、把$H$作为正确率$0.8$的独立因子，观测“规则组报甲、$H$报乙”后，乙类质量为$0.5\times0.3\times0.8=0.12$，甲类质量为$0.5\times0.7\times0.2=0.07$，归一化后$\mathbb{P}(Y=\text{乙}\mid\bm{\lambda})=0.632$。三种结果不同，来自三种依赖假设；这些数值只用于复算，不能替代依赖结构和能力参数的估计。

处理相关性有三种常见粒度。完全重复的来源可以合并，只保留一份票；共享机制清楚但来源仍有独有信息时，可以先按来源组分配总权重；数据与模型允许时，还可以显式加入成对或高阶依赖因子。合并过度会丢失来源差异，依赖结构过细则增加参数量和推断成本。弱监督中的依赖结构学习同样需要稀疏性、样本量和模型条件，不能把“可学习相关性”理解为任意依赖都能从无标签数据恢复\scite{varma2019dependencies}

\section{多源知识的对齐}
\label{sec:knowledge-fusion-alignment}

来源模型回答“谁在怎样作答”，对齐则回答“这些答案是否在讨论同一件事”。一条知识记录至少应保留对象定位、知识内容、来源、任务规范和时间版本。多源对齐的输出不是强制的一一对应，而是对应映射及其状态：确定对应、部分对应、多个候选、待复核或无法转换。只有确定落入共同答案空间的记录才能直接进入投票或概率标签模型。

\subsection{知识对象的对齐}
\label{sec:knowledge-fusion-object-alignment}

\term{对象对齐}（object alignment）确定不同记录指向哪个样本、实体、文本跨度、图像区域或时间区间。对象标识完全相同是最强证据；没有共享标识时，可以比较边界、位置、时间和上下文。对关系知识，还要同时对齐关系两端及方向，因为“甲收购乙”与“乙收购甲”使用相同实体却不是同一命题。

实际匹配通常从精确标识开始，再处理边界或时间范围重叠，最后才在多个候选中使用内容相似度。每一步都应保留原始坐标、匹配依据和置信状态。一条整段标注可能对应多条局部标注，一条时间区间也可能只与另一来源部分重叠；这些一对多和部分对应不能为了便于计票被强行压成一对一。扩大候选搜索范围会提高找回潜在对应的机会，也会增加同名对象错配和计算成本。

图~\ref{fig:knowledge-fusion-object-alignment}使用同一句短文本说明对象错配与标签冲突的区别。左侧两个标签落在不同实体上，虚线把它们配对后产生的“冲突”没有意义；右侧两个来源都标注“李明”，类别才真正发生分歧。实体关系的箭头另行保留端点与方向，不与实体类别标签混合。

\begin{figure*}[htbp]
  \centering
  \begin{tikzpicture}[
  x=1mm,y=1mm,
  every node/.style={font=\figurefont\small,text=BookInk},
  token/.style={book node,draw=BookRule,fill=white,minimum height=8mm,inner sep=2pt},
  span a/.style={draw=BookBlue,fill=FigureInput!10,line width=0.9pt,minimum height=7mm,inner sep=1.5pt},
  span b/.style={draw=BookAmber,fill=BookAmber!10,line width=0.9pt,minimum height=7mm,inner sep=1.5pt}
]
  \node[font=\figurefont\bfseries] at (36,28) {对象错配造成伪冲突};
  \node[token,minimum width=14mm] (ll1) at (9,13) {林舟};
  \node[token,minimum width=9mm] (join1) at (24,13) {在};
  \node[token,minimum width=13mm] (sh1) at (37,13) {上海};
  \node[token,minimum width=9mm] (verb1) at (49,13) {加入};
  \node[token,minimum width=22mm] (co1) at (66,13) {海岬科技};
  \node[span a,minimum width=20mm] (a1) at (9,-2) {来源A：人物};
  \node[span b,minimum width=26mm] (b1) at (66,-16) {来源B：组织};
  \draw[BookBlue,line width=0.9pt] (ll1.south) -- (a1.north);
  \draw[BookAmber,line width=0.9pt] (co1.south) |- (b1.north);
  \draw[book arrow,draw=BookAmber,dashed] (a1.south east) -- node[below,font=\figurefont\scriptsize,text=BookAmber] {错误对应} (b1.north west);
  \node[font=\figurefont\scriptsize,text=BookMuted,align=center] at (36,-29) {标签落在不同实体上\\不能作为同一对象的类别冲突};

  \draw[BookRule,line width=0.7pt] (80,-35) -- (80,32);

  \node[font=\figurefont\bfseries] at (121,28) {对齐后识别真实分歧};
  \node[token,minimum width=14mm] (ll2) at (94,13) {林舟};
  \node[token,minimum width=9mm] (join2) at (109,13) {在};
  \node[token,minimum width=13mm] (sh2) at (122,13) {上海};
  \node[token,minimum width=9mm] (verb2) at (134,13) {加入};
  \node[token,minimum width=22mm] (co2) at (151,13) {海岬科技};
  \node[span a,minimum width=20mm] (a2) at (94,-2) {来源A：人物};
  \node[span b,minimum width=20mm] (b2) at (94,-16) {来源B：组织};
  \draw[BookBlue,line width=0.9pt] (ll2.south) -- (a2.north);
  \draw[BookAmber,line width=0.9pt] (a2.south) -- (b2.north);
  \draw[BookTeal,line width=1pt,-{Stealth[length=2mm]}] (ll2.south east) to[out=-25,in=-155] node[below,font=\figurefont\scriptsize,text=BookTeal] {加入：主体$\to$组织} (co2.south west);
  \node[font=\figurefont\scriptsize,text=BookMuted,align=center] at (128,-29) {同一跨度上的“人物/组织”\\才是需要融合的类别分歧};
\end{tikzpicture}

  \caption{对象错配与真实标签冲突。上、下两行分别表示两个来源的标注跨度；左面板的虚线把不同实体错误配成同一对象，右面板先用实线对齐“李明”，再显示“人物”与“组织”的类别分歧。关系箭头单独标明两端和方向。}
  \label{fig:knowledge-fusion-object-alignment}
\end{figure*}
\FloatBarrier

别名相同不保证对象相同，同名也不保证可以合并；同一个实体在不同时间的状态还可能发生变化。阈值匹配应同时检查漏配和错配，并为边界样本安排复核。对象对齐成功率只衡量映射过程，不能直接当作知识正确率：两条记录可以准确对应同一对象，却一起给出错误标签。

\subsection{知识语义的对齐}
\label{sec:knowledge-fusion-semantic-alignment}

\term{语义对齐}（semantic alignment）确定来源输出是否回答同一个问题，并建立类别定义之间的映射。标签名称只是线索，判定规范、任务条件和有效版本才决定含义。“正类”可以表示文本提到了某项内容，也可以表示该内容在现实中成立；名称相同而命题不同，增加投票次数不能消除差异。

设来源$s$的类别集合为$Y_s$，目标共同空间为$Y^\star$。语义映射$h_s:Y_s\to 2^{Y^\star}$可以把来源类别映到一个目标类别或候选集合。同义类别可作一对一映射；来源类别是上位概念时，只能约束若干细类候选；两个定义交叠但互不包含时，应保留交叠条件和无法映射部分。还要记录$h_s$是否可逆、是否依赖上下文和规范版本，避免映射更新后无法追溯旧结果。

语义映射先作用于类别表或候选集合，再进入融合。若来源A的“相关”表示“文本出现关键词”，来源B的“相关”表示“文本证实命题”，二者不能直接映到同一个二元标签；前者可以作为后者的一项证据或候选约束。无预设语义的簇编号和表示向量更不能因编号、维度相同就逐项对应，它们的含义需要通过生成机制、锚点样本或下游目标建立。

任务规范会随领域和时间修订。新版本把一个旧类别拆成两个细类时，旧记录通常只能映射到候选集合，不能反推出某个新细类。对齐系统应保留原始类别、规范版本和转换链，并在任务迁移时重新检查可比较范围；表面统一的编码若掩盖语义差异，后续概率模型再精细也只是在错误命题上计算。

\subsection{知识粒度与形式的对齐}
\label{sec:knowledge-fusion-granularity-form-alignment}

\term{粒度对齐}（granularity alignment）处理类别层级和对象范围的粗细差异，\term{形式对齐}（form alignment）处理硬标签、概率、分数、候选集合与异构表示之间的表达差异。二者都需要写明转换条件与信息损失。细类概率只有在子类互斥且穷尽父类时才能相加成父类概率；已知父类标签时却不能凭空把概率分配给某个子类。

局部知识到整体知识同样依赖任务规则。一段文本中存在一个正向证据，能否推出整篇文档为正类，要由任务规定的聚合语义决定。概率硬化需要阈值或决策损失，任意模型分数则必须先说明尺度和校准含义。来源没有被调用、调用失败和主动拒答代表不同生成条件，应分别保留，不能用同一个空值进入模型。

表~\ref{tab:knowledge-fusion-form-conversions}汇总常见转换。不可逆转换并非禁止使用，而是要求保留原始记录与转换配置，使后续能够判断不确定性在哪一步丢失。无法合理进入共同答案空间的知识不应勉强投票，第~\ref{sec:knowledge-fusion-heterogeneous-use}节会把它们组织成异构训练目标。

\begin{table*}[htbp]
  \centering
  \small
  \begin{tabularx}{\linewidth}{@{}>{\raggedright\arraybackslash}p{25mm}>{\raggedright\arraybackslash}X>{\raggedright\arraybackslash}X>{\raggedright\arraybackslash}X@{}}
    \toprule
    \textbf{转换} & \textbf{允许操作与条件} & \textbf{保留的信息} & \textbf{无法恢复或需额外依据的信息} \\
    \midrule
    细类到粗类 & 互斥且穷尽的子类可按父类映射合并，概率可在父类内求和 & 父类归属与父类总概率 & 原细类身份及细类间概率分配 \\
    粗类到候选细类 & 将父类映为其全部合法子类候选 & 真值位于哪个父类 & 具体子类；缩小集合需新增证据 \\
    概率到硬标签 & 按阈值、最大后验或明确损失作决策 & 最终行动或训练类别 & 类别间支持程度与冲突强度 \\
    局部到整体 & 仅在任务给出存在、计数或聚合规则时转换 & 规则定义的整体结论 & 未被聚合规则表达的局部位置与数量 \\
    异构表示到共同空间 & 需要配对锚点、已知生成关系或可验证映射 & 映射明确覆盖的共同结构 & 仅凭维度相同不能恢复语义对应 \\
    \bottomrule
  \end{tabularx}
  \caption{知识形式转换的条件与信息损失。每次转换都应保留原始记录、任务条件和转换版本；“需额外依据”表示当前知识只能形成约束或候选，不能直接生成更细答案。}
  \label{tab:knowledge-fusion-form-conversions}
\end{table*}
\FloatBarrier

\section{融合知识的推断}
\label{sec:knowledge-fusion-inference}

完成来源建模与知识对齐后，才能比较不同融合方法。下面固定共同答案空间$Y$和表~\ref{tab:knowledge-fusion-source-records}的来源输出，观察方法改变时哪些假设随之改变。投票把来源视为等价，显式加权区分贡献，概率模型进一步描述潜在标签如何产生来源输出，结构约束则限制多个局部答案能否组成合法整体。这些机制可以组合，不是彼此排斥的系统类别。

\subsection{基于投票的知识融合}
\label{sec:knowledge-fusion-voting}

\term{多数投票}（majority voting）把每个有效来源视为一票。对样本$i$和类别$k$，票数为
\begin{equation}
  v_i(k)
  =\sum_{s=1}^{m}o_{i,s}\mathbb{I}\{\lambda_{i,s}=k\},
  \qquad
  \widehat y_i\in\argmax_{k\in Y}v_i(k).
  \label{eq:knowledge-fusion-majority-vote}
\end{equation}
多类别任务中的相对多数只要求获胜类别票数最多，绝对多数还要求它超过有效票的一半，两种规则不能混用。多标签任务可以逐标签计票，但逐项通过的组合仍可能违反互斥、层级或序列条件。关系知识也只有在对齐成同一关系命题及方向后才能计票。

投票必须明确平票、全部缺失和最低票数规则。表~\ref{tab:knowledge-fusion-comparison}采用相对多数：唯一最高票时输出该类，平票或零张有效票时保留未决。弃权与失败不进入分母，票数比例只表示来源输出比例，不是经过校准的真实正确概率。

\begin{table*}[htbp]
  \centering
  \small
  \begin{tabularx}{\linewidth}{@{}ccccc>{\raggedright\arraybackslash}X@{}}
    \toprule
    \textbf{样本} & \textbf{有效票} & \textbf{甲票/乙票} & \textbf{等权结果} & \textbf{加权结果} & \textbf{关键原因} \\
    \midrule
    $x_1$ & $4$ & $4/0$ & 甲 & 甲 & 全部来源一致 \\
    $x_2$ & $3$ & $2/1$ & 甲 & 乙 & 模型权重高于共享规则组总权重 \\
    $x_3$ & $4$ & $3/1$ & 甲 & 甲 & 人工与规则共同支持甲 \\
    $x_4$ & $1$ & $1/0$ & 甲 & 甲 & 只有模型作答，权重不能补充缺失证据 \\
    $x_5$ & $3$ & $2/1$ & 甲 & 乙 & 两张甲票来自同一规则模板 \\
    $x_6$ & $4$ & $2/2$ & 未决 & 乙 & 人工与模型共同支持乙 \\
    $x_7$ & $0$ & $0/0$ & 未决 & 未决 & 全部来源缺失或弃权 \\
    $x_8$ & $4$ & $2/2$ & 未决 & 甲 & 独立人工与模型共同支持甲 \\
    \bottomrule
  \end{tabularx}
  \caption{同一组候选的融合结果对照。加权列使用教学权重$w_{R_1}=w_{R_2}=0.25$、$w_H=1.2$、$w_M=0.6$，只用于复核运算与展示假设差异，不表示这些权重经过真实审计，也不证明加权方法普遍优于投票。}
  \label{tab:knowledge-fusion-comparison}
\end{table*}
\FloatBarrier

投票的计算成本低、过程容易复核，适合作为融合基线，但它略去了三类信息：来源能力不同，错误方向不同，来源之间可能相关。表中的$x_5$最能说明第三点；复制一条规则就能把一张甲票变成两张，结果却没有增加独立证据。来源数量增加能否改善结果，取决于新增来源的能力、覆盖和条件依赖，不能只看票数。

若直接遍历$n\times m$的稠密标签矩阵，计票时间为$O(nm)$；只存有效输出时，可以写成$O(N_{\mathrm{obs}}+nK)$，其中$N_{\mathrm{obs}}=\sum_{i,s}o_{i,s}$用于累计有效票，$nK$用于在每个样本的$K$个类别中选最大值。逐样本处理只需$O(K)$工作内存，保存全部票数则需$O(nK)$。相关来源去重和语义对齐的成本不包含在这项计票复杂度中。

\subsection{基于加权的知识融合}
\label{sec:knowledge-fusion-weighting}

\term{加权融合}（weighted fusion）为来源贡献分配不同权重。硬标签的加权得分可以写成
\begin{equation}
  s_i(k)
  =\sum_{s=1}^{m}o_{i,s}w_{i,s,k}
    \mathbb{I}\{\lambda_{i,s}=k\},
  \qquad
  \widehat y_i\in\argmax_{k\in Y}s_i(k).
  \label{eq:knowledge-fusion-weighted-vote}
\end{equation}
权重$w_{i,s,k}$可以是来源全局权重，也可以依赖类别$k$、样本区域或时间。它应来自独立审计、已明确的来源模型或预先规定的设计规则，并说明量纲。选择权重所用的审计数据不能同时充当最终质量评价数据。

若来源输出的是类别概率$p_s(k\mid x_i)$，一种线性汇总为
\begin{equation}
  \widetilde p_i(k)
  =\frac{\sum_{s=1}^{m}o_{i,s}w_{i,s}p_s(k\mid x_i)}
  {\sum_{s=1}^{m}o_{i,s}w_{i,s}}.
  \label{eq:knowledge-fusion-weighted-probability}
\end{equation}
只有当各来源概率的类别语义、尺度和校准含义可比，且权重非负、分母为正时，这个表达才直接成立。任意模型分数不能在没有转换依据时当作概率平均；即使式~\eqref{eq:knowledge-fusion-weighted-probability}归一化为一，也不自动获得关于真实频率的校准保证。

表~\ref{tab:knowledge-fusion-comparison}的教学权重让人工$H$承担更大贡献，并降低共享模板规则的单票权重。以$x_5$为例，甲类得分为$0.25+0.25=0.5$，乙类得分为$1.2$，所以加权结果为乙。该结果来自显式权重假设，并不说明人工标签必然正确。

高置信也不等于高可靠。考虑一个参考答案为乙的教学样本：存在系统性甲类偏差的来源$C$输出$p_C=(0.99,0.01)$，独立来源$D$输出$p_D=(0.25,0.75)$，向量分量依次对应甲、乙。若把两者各自的最大置信度$0.99$和$0.75$直接当来源权重，甲、乙的未归一化得分分别为$0.99\times0.99+0.75\times0.25=1.1676$和$0.99\times0.01+0.75\times0.75=0.5724$，错误的甲类仍获胜。若独立审计支持可靠性权重$w_C=0.2,w_D=0.8$，两类得分则变为$0.398$与$0.602$，结果转为乙。这个对照说明置信度描述来源对当前输出的集中程度，审计权重描述来源在目标范围内的经验可靠性；后二者仍须分别检查概率校准和审计代表性。

权重复杂化会带来估计和维护成本。类别相关权重需要每个类别上的证据，区域相关权重还要处理区域覆盖；负权重会把来源解释成反向证据，必须有明确模型支持。实践中应检查去掉单个来源、合并相关来源、更新审计数据或扰动权重后结果是否稳定，并防止少数权重过度集中。加权解决的是“来源贡献不同”，仍没有显式描述观察怎样由潜在真值产生。

\subsection{基于概率模型与约束的知识融合}
\label{sec:knowledge-fusion-probabilistic-constrained}

概率模型把真实标签作为未观测变量，并用来源参数解释观察到的标签；结构约束则规定多个局部标签怎样组成可接受的整体。前者回答每个候选得到多少概率支持，后者回答哪些候选组合允许出现。联合系统可以先由概率模型形成局部得分，再在硬约束或软约束下寻找整体赋值。

\subsubsection{潜在真实标签的概率推断}
\label{sec:knowledge-fusion-latent-label-inference}

\term{Dawid--Skene模型}把每个样本的未知真实类别$Y_i$设为潜变量，用每个来源的混淆矩阵描述观察误差。设类别先验为$\pi_k=\mathbb{P}(Y_i=k)$，来源$s$的混淆参数为式~\eqref{eq:knowledge-fusion-confusion-matrix}中的$\theta^{(s)}_{k,\ell}$，实际对样本$i$作答的来源集合为$S_i$。标准模型假设样本独立同分布、来源混淆矩阵在样本间稳定，并且给定$Y_i$后来源输出条件独立；若缺失机制不另行建模，还要把已观察的来源集合$S_i$视为固定，或假设缺失机制对当前似然可忽略。于是完整数据联合分布与观测似然分别为
\begin{equation}
  p(Y_i=k,\bm{\lambda}_i)
  =\pi_k\prod_{s\in S_i}\theta^{(s)}_{k,\lambda_{i,s}},
  \qquad
  p(\bm{\lambda}_i)
  =\sum_{k=1}^{K}\pi_k
    \prod_{s\in S_i}\theta^{(s)}_{k,\lambda_{i,s}}.
  \label{eq:knowledge-fusion-dawid-skene}
\end{equation}
这组假设明确了模型边界：来源共享错误时条件独立失效，来源只在某些难度区域作答时固定混淆与可忽略缺失可能失效，未在任何重叠中出现的类别--来源组合也没有直接估计证据。Dawid与Skene用EM算法联合估计观察者误差率和潜在类别\scite{dawid1979observer}

EM的E步在当前参数$\bm{\theta}^{(t)}$下计算类别责任度
\begin{equation}
  q_{i,k}^{(t)}
  =
  \frac{\pi_k^{(t)}
    \prod_{s\in S_i}
    \left(\theta^{(s)}_{k,\lambda_{i,s}}\right)^{(t)}}
  {\sum_{r=1}^{K}\pi_r^{(t)}
    \prod_{s\in S_i}
    \left(\theta^{(s)}_{r,\lambda_{i,s}}\right)^{(t)}}.
  \label{eq:knowledge-fusion-dawid-skene-e-step}
\end{equation}
M步把未知真值计数替换为责任度加权计数：
\begin{equation}
  \pi_k^{(t+1)}
  =\frac{1}{n}\sum_{i=1}^{n}q_{i,k}^{(t)},
  \qquad
  \left(\theta^{(s)}_{k,\ell}\right)^{(t+1)}
  =
  \frac{\sum_{i:s\in S_i}
    q_{i,k}^{(t)}\mathbb{I}\{\lambda_{i,s}=\ell\}}
  {\sum_{i:s\in S_i}q_{i,k}^{(t)}}.
  \label{eq:knowledge-fusion-dawid-skene-m-step}
\end{equation}
分母为零表示当前数据与参数没有为该行提供有效质量，不能照抄更新值；可以保留未估计状态或加入明确先验平滑。初始化可用多数投票、少量审计标签或接近对角的混淆矩阵。常见停止条件是观测对数似然或参数变化低于阈值，也应设置最大迭代次数。

下面完整手算一轮二分类EM。设三个来源$s=1,2,3$的初始混淆均为对称误差，正确率依次为$0.8,0.7,0.6$，类别先验各为$0.5$；三个样本的观察依次为$(\text{甲},\text{甲},\text{甲})$、$(\text{甲},\text{乙},\text{乙})$和$(\text{乙},\text{乙},\text{甲})$。第一个样本在甲、乙下的未归一化质量分别为$0.5\times0.8\times0.7\times0.6=0.168$与$0.5\times0.2\times0.3\times0.4=0.012$；第二个样本对应$0.048$与$0.042$；第三个样本对应$0.018$与$0.112$。归一化后，三个甲类责任度为$0.9333,0.5333,0.1385$，乙类责任度为各自的补数。

M步得到$(\pi_{\text{甲}}^{(1)},\pi_{\text{乙}}^{(1)})=(0.5350,0.4650)$。为展示混淆更新，甲类软计数总量为$1.6051$，乙类软计数总量为$1.3949$；例如来源1在前两个样本输出甲，因此$\left(\theta^{(1)}_{\text{甲},\text{甲}}\right)^{(1)}=(0.9333+0.5333)/1.6051=0.9137$，它在这两个样本上的乙类软计数为$0.0667+0.4667$，故$\left(\theta^{(1)}_{\text{乙},\text{甲}}\right)^{(1)}=0.5334/1.3949=0.3824$。为紧凑展示，记来源$s$在第1次M步后的矩阵为$\bm{\Theta}_s^{(1)}=\left(\left(\theta^{(s)}_{k,\ell}\right)^{(1)}\right)_{k,\ell}$。按“行是真实类别、列是来源输出甲/乙”的约定，三个更新矩阵为
\[
  \begin{aligned}
    \bm{\Theta}_1^{(1)}
    &=
    \begin{pmatrix}
      0.9137 & 0.0863 \\
      0.3824 & 0.6176
    \end{pmatrix},\\
    \bm{\Theta}_2^{(1)}
    &=
    \begin{pmatrix}
      0.5815 & 0.4185 \\
      0.0478 & 0.9522
    \end{pmatrix},\\
    \bm{\Theta}_3^{(1)}
    &=
    \begin{pmatrix}
      0.6677 & 0.3323 \\
      0.6654 & 0.3346
    \end{pmatrix}.
  \end{aligned}
\]
每一行都归一化为一，至此先验和三个来源的两行混淆参数都完成了一次更新。这里没有加入平滑；三个样本只够复算步骤，不足以稳定估计模型。

图~\ref{fig:knowledge-fusion-latent-label-em}把固定观察、当前参数和潜在标签后验分开。标准EM在精确完成E步和M步时不会降低当前模型的观测似然，但似然通常非凸，停止可能只得到局部最优或鞍点附近结果。数值收敛只说明更新趋于稳定，不说明参数可识别，更不说明潜在答案与现实真值一致；三者需要分别检查。

\begin{figure}[htbp]
  \centering
  \begin{tikzpicture}[
  x=1mm,y=1mm,
  every node/.style={font=\figurefont\small,text=BookInk},
  data/.style={book node,draw=BookBlue,fill=FigureInput!10,line width=0.9pt,minimum width=42mm,minimum height=11mm,inner sep=2pt},
  state/.style={book node,draw=BookTeal,fill=FigureModel!10,line width=0.9pt,minimum width=45mm,minimum height=11mm,inner sep=2pt},
  check/.style={book node,draw=BookAmber,fill=BookAmber!10,line width=0.9pt,minimum width=34mm,minimum height=10mm,inner sep=2pt},
  output/.style={book node,draw=BookInk,fill=BookPaper,line width=0.9pt,minimum width=39mm,minimum height=10mm,inner sep=2pt}
]
  \node[data] (matrix) at (0,34) {固定观测$\bm{\Lambda},\bm{O}$};
  \node[state] (init) at (0,17) {初始化$\bm{\pi}^{(0)},\bm{\Theta}^{(0)}$};
  \node[state] (estep) at (0,0) {E步：计算$q_{i,k}^{(t)}$};
  \node[state] (mstep) at (0,-17) {M步：更新$\bm{\pi}^{(t+1)},\bm{\Theta}^{(t+1)}$};
  \node[check] (stop) at (0,-34) {似然或参数已稳定？};
  \node[output] (result) at (0,-51) {后验分布与来源参数};

  \draw[book arrow,knowledge arrow,draw=BookBlue,line width=1pt] (matrix) -- (init);
  \draw[book arrow,knowledge arrow,draw=BookTeal,line width=1pt] (init) -- (estep);
  \draw[book arrow,knowledge arrow,draw=BookTeal,line width=1pt] (estep) -- node[right,font=\figurefont\scriptsize] {软类别计数} (mstep);
  \draw[book arrow,knowledge arrow,draw=BookTeal,line width=1pt] (mstep) -- (stop);
  \draw[book arrow,knowledge arrow,draw=BookInk,line width=1pt] (stop) -- node[right,font=\figurefont\scriptsize] {是} (result);
  \draw[book arrow,knowledge arrow,draw=BookBlue] (matrix.east) -- ++(16,0) |- node[pos=0.22,right,font=\figurefont\scriptsize,text=BookBlue] {每轮读取} (estep.east);
  \draw[book arrow,knowledge arrow,draw=BookAmber,dashed] (stop.west) -- ++(-16,0) |- node[pos=0.18,left,font=\figurefont\scriptsize,text=BookAmber,align=right] {否：携带\\新参数} (estep.west);

  \node[anchor=west,font=\figurefont\scriptsize,text=BookMuted,align=left] at (32,-42) {分别检查：\\可识别性\\数值收敛\\答案正确性};
\end{tikzpicture}

  \caption{来源参数与潜在标签的交替估计。观测标签矩阵在迭代中保持不变；E步用当前来源参数计算类别后验，M步用后验软计数更新类别先验和来源参数，再检查停止条件。实线表示数据与计算流，虚线表示更新后的参数回流。}
  \label{fig:knowledge-fusion-latent-label-em}
\end{figure}
\FloatBarrier

Dawid--Skene模型至少存在潜在类别置换对称性：同时置换$\pi$的元素、所有混淆矩阵的行和潜在标签名称，$p(\bm{\Lambda})$不变。参考锚点或方向假设负责命名类别；来源重叠不足、混淆矩阵退化或相关来源被误当独立，还会造成额外不可识别或弱识别。第~\ref{sec:knowledge-fusion-relational-reliability}节的秩条件说明，恢复能力来自足够多、足够不同的条件分布，而不是来自EM本身。

\term{标签模型}（label model）是对弱来源输出建立的生成模型家族，不是Dawid--Skene的简单别名。令标注函数输出$\lambda_{i,s}\in Y\cup\{\bot\}$，一个通用对数线性形式为
\begin{equation}
  p_{\bm{\theta}}(Y_i,\bm{\lambda}_i)
  =\frac{1}{Z_{\bm{\theta}}}
  \exp\left(
    \sum_{r=1}^{R}\theta_r
    \phi_r(Y_i,\bm{\lambda}_i)
  \right).
  \label{eq:knowledge-fusion-label-model}
\end{equation}
因子$\phi_r$可以描述类别先验、来源准确性、类别偏好、覆盖或弃权倾向，以及特定来源对的相关性。最简单标签模型只保留类别先验与条件独立准确性因子；更完整的模型可加入已知或学得的依赖结构。Data Programming把规则、知识库和模型等写成可弃权的标注函数，再从它们的重叠与冲突中估计生成模型参数\scite{ratner2016data-programming} Snorkel进一步把这种标签模型产生的概率标签交给下游判别模型\scite{ratner2017}

式~\eqref{eq:knowledge-fusion-label-model}的标准用法还假设各样本在固定参数下独立抽取，同一来源的准确性、覆盖与依赖参数在所建模范围内保持稳定，并且$\bot$明确表示标注函数运行后选择不投票。无标签拟合实际使用边缘分布$p_{\bm{\theta}}(\bm{\lambda}_i)=\sum_k p_{\bm{\theta}}(Y_i=k,\bm{\lambda}_i)$；基础标签模型只观察输出矩阵，不会自动利用原始输入$\bm{x}_i$判断某个样本更难。若来源质量依赖样本或时间，应分组、加入相应因子或改用实例相关模型，而不能把变化都吸收到一个全局准确性参数中。

标签模型与Dawid--Skene共享潜在类别、来源误差和无参考估计的核心思想，但参数化重点不同。Dawid--Skene通常为每个来源使用完整类别混淆矩阵；程序化弱监督的标签模型常用较少的准确性因子，显式允许弃权，并可加入来源依赖。完整混淆更能表达错误方向，也需要更多重叠证据；较粗准确性因子减少参数，却加入了对称误差或共享结构等额外限制。

标签模型的可识别性同样不是无条件保证。除了类别置换，对数线性参数还可能存在不同参数给出同一分布的冗余表示；条件独立版本需要足够重叠和有效独立来源，依赖版本还要求因子结构正确且数据足以区分依赖参数。Data Programming给出的无标签参数恢复结论成立于论文规定的模型与条件，不能推广成“任意几条规则都能恢复真值”。来源高度相关、平均方向未知或参数量超过有效证据时，应增加锚点、简化模型、合并来源或保留未决，而不是把一个收敛的后验当作验证结果。

\subsubsection{结构约束下的联合推断}
\label{sec:knowledge-fusion-structured-inference}

逐样本或逐位置融合得到的是局部证据，但任务常要求答案之间满足结构。类别层级规定子类出现时父类也应成立，互斥关系禁止两个标签同时选择，BIO（beginning--inside--outside）序列标注则不允许$I$标签在没有相应$B$标签的情况下开始。若分别取每个位置的最高分，组合可能不合法，因此需要\term{联合推断}（joint inference）：一次确定一组相互耦合的标签。

前面的基于规则获取关注约束怎样从已有知识产生候选知识，本节则把多源局部得分视为已经给定，关注约束怎样从候选组合中确定整体结果。两处都可以使用同一条业务规则，但输入、输出和评价对象不同，不能因为都出现约束就把两个过程合并。

设局部融合为变量$z_i\in Y$提供得分$a_i(z_i)$，硬约束定义由可行标签向量组成的集合$\bm{Z}_{\mathrm{feas}}$，软约束$r$的违反代价为$\rho_r(\bm{z})\geq0$、权重为$\eta_r\geq0$。统一目标可以写成
\begin{equation}
  \widehat{\bm{z}}
  \in\argmax_{\bm{z}\in\bm{Z}_{\mathrm{feas}}}
  \left\{
    \sum_i a_i(z_i)
    -\sum_r\eta_r\rho_r(\bm{z})
  \right\}.
  \label{eq:knowledge-fusion-structured-objective}
\end{equation}
硬约束通过$\bm{Z}_{\mathrm{feas}}$排除候选，软约束只改变候选偏好。若概率模型给出$q_{i,k}$，局部得分可以取$\log q_{i,k}$；若输入只是未校准加权分数，则目标值也只能解释为当前评分体系下的偏好。

考虑三个词元的BIO序列。令第三个位置只允许$O$，其余有限局部分数为
\[
  \begin{array}{c|ccc}
    & O & B\text{-ORG} & I\text{-ORG} \\
    \hline
    1 & 1.2 & 0.8 & -\infty \\
    2 & 0.5 & 0.6 & 1.4 \\
    3 & 0.7 & -\infty & -\infty
  \end{array}.
\]
逐位置最大化得到$(O,I\text{-ORG},O)$，局部分数为$3.3$，但$I\text{-ORG}$不能跟在$O$之后。有限分数下的全部合法候选及总分为：$(O,O,O)$得$2.4$，$(O,B\text{-ORG},O)$得$2.5$，$(B\text{-ORG},O,O)$得$2.0$，$(B\text{-ORG},B\text{-ORG},O)$得$2.1$，$(B\text{-ORG},I\text{-ORG},O)$得$2.9$。硬约束因此选择最后一个候选，而不是简单删除中间标签。

若把“$O$后不能接$I\text{-ORG}$”改为一次违反记$1$的软代价，原局部最优的目标值为$3.3-\eta$，其余上述合法候选的代价为零。当$\eta=0.2$时，非法序列仍以$3.1$胜过$2.9$；当$\eta=0.5$时，它降为$2.8$，联合最优改为合法的$(B\text{-ORG},I\text{-ORG},O)$。这次可枚举求解同时说明：硬约束直接改变候选集合，软约束则用可核对的罚值控制局部证据与结构偏离之间的取舍。互斥分类中也类似：两个局部正标签同时得分最高时，可以选择得分损失较小的一项，或在约束可能有误时把互斥写成软代价。

约束本身也是知识来源，可能过时、冲突或只在部分对象上适用。求解失败时要区分三种情况：硬约束确实没有可行赋值，模型把不适用约束施加给了当前样本，或求解器在资源限制下没有找到解。直接删除冲突会丢失来源，放松约束需要记录违反程度，近似求解则要说明未保证的是最优性、可行性还是二者。变量耦合越密，联合搜索成本越高；链式序列可用动态规划，通用离散约束可能需要整数规划、消息传递或近似搜索。

\section{融合结果的表达与使用}
\label{sec:knowledge-fusion-result-use}

融合推断得到的可以是票数、分数、后验或整体赋值，后续系统还要决定怎样使用它。把“模型算出了什么”与“系统采取什么行动”分开，才能解释同一份证据为何可以形成硬标签、软训练目标、候选集合或弃权。无论选择哪种表达，都应保留来源摘要、对齐映射、融合配置和结果版本。

\subsection{单一标签与概率分布的使用}
\label{sec:knowledge-fusion-label-probability-use}

若当前任务必须作出一个离散行动，可以选择后验最大的类别；若不同错误代价不同，则应最小化后验期望损失。设融合后类别分布为$q_i(k)$，行动集合为$A$，把真实类别$k$时采取行动$a$的损失记为$L(a,k)$，贝叶斯决策为
\begin{equation}
  \widehat a_i
  \in\argmin_{a\in A}
  \sum_{k=1}^{K}L(a,k)q_i(k).
  \label{eq:knowledge-fusion-bayes-decision}
\end{equation}
最大后验只是零一损失下的特殊情形。漏判代价高于误报代价时，即使两个样本后验相同，采用不同业务损失也可能得到不同阈值；这属于使用决策，不应改写成融合模型“更确信”某一类。

若结果用于训练，保留$q_i$可以形成软标签交叉熵$-\sum_k q_i(k)\log p_{\bm{w}}(k\mid\bm{x}_i)$，让各类别支持程度都进入目标。硬化为$\argmax_k q_i(k)$会舍弃剩余质量和冲突强度，但接口更简单、存储更少。若$q_i$来自未校准加权分数或错设来源模型，只能称为当前模型下的支持分布，不能自动解释为真实频率。

概率标签还需要版本管理。来源新增、混淆更新或语义映射修订后，同一样本的分布可能改变；保存原始分布、决策阈值、损失矩阵和模型版本，才能在审计时区分“推断结果变了”与“使用规则变了”。训练系统若只接收硬标签，可以在接口处硬化，但不应删除上游概率记录。

\subsection{候选集合与弃权的使用}
\label{sec:knowledge-fusion-set-abstention-use}

证据只能排除部分答案时，可以输出候选集合$C_i\subseteq Y$。一种规则是按$q_i(k)$从高到低加入类别，直到累计质量达到阈值$\tau$；也可以保留所有$q_i(k)\geq\gamma$的类别。空集合表示规则设置不合适或没有候选满足阈值，单元素集合退化为硬标签，过大的集合则可能对下游没有区分价值。候选集合是否覆盖真值仍需要独立校准与假设，不能因名称叫“置信集合”就获得保证。

\term{弃权}（abstention）表示融合系统在当前用途要求下不输出可用决策。来源自身弃权发生在输入端，融合弃权发生在综合证据后，二者必须分别记录。融合系统可以在$\max_k q_i(k)<\eta$、有效独立来源少于阈值或对齐状态未确认时弃权，并触发人工复核、追加来源或保留未决。弃权不是负类，也不应被永久排除在质量统计之外。

图~\ref{fig:knowledge-fusion-output-forms}让同一教学分布$q=(0.62,0.28,0.10)$分别进入四种使用规则。最大后验给出硬标签甲；软目标保留全部分量；累计质量阈值$\tau=0.90$得到候选集合$\{\text{甲},\text{乙}\}$；高风险用途若要求$\max q\geq0.70$才接收，则当前样本弃权。这些分支使用不同规则，不表示系统会同时作出互斥行动。

\begin{figure*}[htbp]
  \centering
  \begin{tikzpicture}[
  x=1mm,y=1mm,
  every node/.style={font=\figurefont\small,text=BookInk},
  evidence/.style={book node,draw=BookInk,fill=BookPaper,line width=1pt,minimum width=61mm,minimum height=12mm,inner sep=2pt},
  result/.style={book node,minimum width=33mm,minimum height=18mm,inner sep=2pt,text width=30mm,align=center}
]
  \node[evidence] (q) at (75,19) {教学支持分布（山茶、茶梅、月季）\\$q=(0.62,\ 0.28,\ 0.10)$};

  \node[result,draw=BookBlue,fill=FigureInput!9,line width=0.9pt] (hard) at (12,-14) {硬标签\\山茶};
  \node[result,draw=BookTeal,fill=FigureModel!9,line width=0.9pt] (soft) at (54,-14) {概率目标\\$(0.62,0.28,0.10)$};
  \node[result,draw=BookAmber,fill=BookAmber!9,line width=0.9pt,dashed] (set) at (96,-14) {候选集合\\$\{\text{山茶},\text{茶梅}\}$};
  \node[result,draw=BookMuted,fill=BookPaper,line width=0.9pt,dotted] (abstain) at (138,-14) {暂缓使用\\弃权};

  \draw[book arrow,knowledge arrow,draw=BookBlue,line width=0.9pt] (q.south west) -- ++(0,-4) -| node[pos=0.74,left,font=\figurefont\scriptsize,text=BookBlue] {$\argmax$} (hard.north);
  \draw[book arrow,knowledge arrow,draw=BookTeal,line width=0.9pt] ([xshift=15mm]q.south west) -- ++(0,-6) -| node[pos=0.76,left,font=\figurefont\scriptsize,text=BookTeal] {完整保留} (soft.north);
  \draw[book arrow,knowledge arrow,draw=BookAmber,dashed,line width=0.9pt] ([xshift=-15mm]q.south east) -- ++(0,-6) -| node[pos=0.76,right,font=\figurefont\scriptsize,text=BookAmber] {累计权重$\geq0.90$} (set.north);
  \draw[book arrow,knowledge arrow,draw=BookMuted,dotted,line width=0.9pt] (q.south east) -- ++(0,-4) -| node[pos=0.74,right,font=\figurefont\scriptsize,text=BookMuted] {$\max q<0.70$} (abstain.north);

  \node[font=\figurefont\scriptsize,text=BookMuted] at (12,-31) {离散决策};
  \node[font=\figurefont\scriptsize,text=BookMuted] at (54,-31) {软标签训练};
  \node[font=\figurefont\scriptsize,text=BookMuted] at (96,-31) {部分排除};
  \node[font=\figurefont\scriptsize,text=BookMuted] at (138,-31) {触发复核};
\end{tikzpicture}

  \caption{同一融合证据的不同输出方式。三类别支持分布为人工教学设定；四个分支分别采用最大后验、完整概率、累计质量阈值和最低接收阈值，展示硬化、保留不确定性、保留候选与暂缓使用的差异，不赋予该分布额外校准或覆盖保证。}
  \label{fig:knowledge-fusion-output-forms}
\end{figure*}
\FloatBarrier

提高接收阈值通常会降低覆盖，已接收样本的错误却不保证在每个区域都同步下降。若困难区域更常被弃权，总体已接收准确率可能改善，同时系统对困难对象几乎不提供服务。因此应同时报告总体与分组覆盖、已接收错误和候选集合大小，并用独立数据选择阈值。下游训练若接受候选集合，还要明确集合内类别怎样分配损失，不能把“属于集合”误写成集合内每个类别都正确。

一个人工计数可以说明这种构成变化。假设简单区和困难区各有$100$个样本；阈值为$0.60$时分别接收$90$个和$60$个，其中错误$9$个和$18$个，总体已接收错误率为$27/150=18\%$。把阈值提高到$0.80$后分别接收$70$个和$20$个，其中错误$4$个和$5$个，总体错误率降为$9/90=10\%$，但困难区覆盖从$60\%$降到$20\%$，其已接收错误率仍为$25\%$。这些数字只是教学设定；它们说明总体指标改善可以主要来自少接收困难样本，不能据此推断每个区域都变得更可靠。

\subsection{异构知识的联合使用}
\label{sec:knowledge-fusion-heterogeneous-use}

任务类别、样本对关系、局部标签和表示恢复目标描述的不是同一随机变量，不能先编码成相同维度再平均。\term{异构知识的联合使用}是保留每类知识的监督对象，让它们通过共享模型、多个损失项或结构约束共同影响训练。输出通常是联合训练安排，而不是一个统一融合标签。

设共享模型参数为$\bm{w}$，知识类型$r$的适用样本掩码为$m_{i,r}\in\{0,1\}$，单项损失为$\ell_r(\bm{w};i)$，权重为$\alpha_r\geq0$。一个基本目标为
\begin{equation}
  \mathcal{L}(\bm{w})
  =
  \sum_{r=1}^{R}
  \alpha_r
  \frac{\sum_{i=1}^{n}m_{i,r}\ell_r(\bm{w};i)}
  {\max\left(1,\sum_{i=1}^{n}m_{i,r}\right)}.
  \label{eq:knowledge-fusion-heterogeneous-objective}
\end{equation}
掩码说明哪些样本具有哪类知识，缺失某类知识时跳过对应项，而不是把目标值填成零。分母按有效样本归一化，仍不能保证不同损失尺度可比；权重$\alpha_r$同时受到目标重要性、梯度尺度、采样频率和计算成本影响。

图~\ref{fig:knowledge-fusion-heterogeneous-training}以图像训练为例。类别标签监督分类头，两个视图的已知对应关系监督表示一致性，局部遮挡内容监督恢复头。三类知识分别经过适用样本掩码和权重后汇入总损失，共享的是表示模型，不是标签语义。不同编码器给出的向量若没有锚点映射，也不能直接逐维相加。

\begin{figure*}[htbp]
  \centering
  \begin{tikzpicture}[
  x=0.8mm,y=1mm,
  every node/.style={font=\figurefont\small,text=BookInk},
  input/.style={book node,draw=BookBlue,fill=FigureInput!10,line width=0.9pt,minimum width=22mm,minimum height=10mm,inner sep=2pt},
  model/.style={book node,draw=BookTeal,fill=FigureModel!10,line width=1pt,minimum width=28mm,minimum height=15mm,inner sep=2pt},
  branch/.style={book node,draw=BookTeal,fill=FigureModel!7,line width=0.8pt,minimum width=27mm,minimum height=9mm,inner sep=2pt},
  knowledge/.style={book node,draw=BookAmber,fill=BookAmber!9,line width=0.8pt,minimum width=28mm,minimum height=9mm,inner sep=2pt},
  loss/.style={book node,draw=BookInk,fill=BookPaper,line width=0.8pt,minimum width=31mm,minimum height=10mm,inner sep=2pt}
]
  \node[input] (image) at (0,0) {图像与增强视图};
  \node[model] (encoder) at (37,0) {共享表示模型\\$f_{\bm{w}}$};
  \draw[book arrow,draw=BookBlue,line width=1pt] (image) -- (encoder);

  \node[branch] (pred) at (78,25) {分类预测};
  \node[branch] (compare) at (78,0) {视图表示比较};
  \node[branch] (recover) at (78,-25) {局部内容重建};
  \node[loss] (l1) at (130,25) {$m_1\alpha_1\mathcal{L}_{\mathrm{cls}}$};
  \node[loss] (l2) at (130,0) {$m_2\alpha_2\mathcal{L}_{\mathrm{rel}}$};
  \node[loss] (l3) at (130,-25) {$m_3\alpha_3\mathcal{L}_{\mathrm{rec}}$};
  \node[knowledge] (label) at (105,43) {外部任务类别标签};
  \node[knowledge] (pair) at (105,13) {外部视图对应关系};
  \node[knowledge] (patch) at (105,-43) {外部局部恢复目标};
  \node[loss,draw=BookTeal,fill=FigureModel!12,line width=1pt,minimum width=20mm] (sum) at (165,0) {总损失};

  \draw[book arrow,draw=BookTeal] (encoder.east) -- ++(3,0) |- (pred.west);
  \draw[book arrow,draw=BookTeal] (encoder.east) -- (compare.west);
  \draw[book arrow,draw=BookTeal] (encoder.east) -- ++(3,0) |- (recover.west);
  \draw[book arrow,draw=BookTeal] (pred) -- node[below,font=\figurefont\scriptsize,text=BookTeal] {模型输出} (l1);
  \draw[book arrow,draw=BookTeal] (compare) -- node[below,font=\figurefont\scriptsize,text=BookTeal] {模型输出} (l2);
  \draw[book arrow,draw=BookTeal] (recover) -- node[above,font=\figurefont\scriptsize,text=BookTeal] {模型输出} (l3);
  \draw[book arrow,draw=BookAmber] (label.south east) |- node[pos=0.72,right,font=\figurefont\scriptsize,text=BookAmber] {监督目标} (l1.north);
  \draw[book arrow,draw=BookAmber] (pair.south east) |- node[pos=0.72,right,font=\figurefont\scriptsize,text=BookAmber] {监督关系} (l2.north);
  \draw[book arrow,draw=BookAmber] (patch.north east) |- node[pos=0.72,right,font=\figurefont\scriptsize,text=BookAmber] {监督目标} (l3.south);
  \draw[book arrow,draw=BookTeal,line width=1pt] (l1.east) -- ++(6,0) |- (sum.north);
  \draw[book arrow,draw=BookTeal,line width=1pt] (l2.east) -- (sum.west);
  \draw[book arrow,draw=BookTeal,line width=1pt] (l3.east) -- ++(6,0) |- (sum.south);

  \node[anchor=north,font=\figurefont\scriptsize,text=BookMuted,align=center] at (130,-52) {掩码$m_r$在损失处选择具有相应知识的样本};
\end{tikzpicture}

  \caption{异构知识形成联合训练约束。图像经共享表示模型进入三个损失分支；任务标签、视图关系和局部恢复目标分别提供不同监督对象，经适用样本掩码与权重后组成总训练目标。共享训练不要求把三类知识压缩成同一类别标签。}
  \label{fig:knowledge-fusion-heterogeneous-training}
\end{figure*}
\FloatBarrier

共享参数会让不同目标相互影响，权重改变的不只是损失数值，还会改变训练资源和梯度方向。需要分别检查各目标的适用数据、尺度、采样比例及单独贡献；梯度长期冲突时，可以调整权重、交替批次、分阶段训练或减少共享层。任何方案都应保留每项知识的独立验证指标，否则总损失下降无法说明哪类知识真正发挥了作用。

\subsection{追加来源与跨来源预算的分配}
\label{sec:knowledge-fusion-source-budget}

融合后仍冲突或未覆盖时，可以追加人工判断、补充规则、调用另一模型或继续收集数据。选择依据不是单次调用价格，而是新增来源能否提供当前缺少的信息。再次调用同一确定模型成本很低，却可能完全复制原证据；一名独立专家价格更高，却可能显著改变候选排序。

把当前信息状态记为$\mathcal{D}$，候选行动集合记为$A$，采取行动$a\in A$后可能得到结果$Z_a$，结果质量用任务效用$U(\mathcal{D})$度量。人工时间、货币费用、计算量和反馈时延具有不同量纲，分别记作$C_j(a)$，可用预算记作$B_j$，$j=1,\ldots,J$。不做单位换算时，应在各自预算约束下最大化期望任务效用增量：
\begin{equation}
  a^\star
  \in\argmax_{a\in A}
  \left\{
    \mathbb{E}_{Z_a\mid\mathcal{D}}
    [U(\mathcal{D}\cup Z_a)]
    -U(\mathcal{D})
  \right\}
  \quad
  \text{s.t.}\quad
  C_j(a)\leq B_j,\ j=1,\ldots,J.
  \label{eq:knowledge-fusion-source-budget}
\end{equation}
式~\eqref{eq:knowledge-fusion-source-budget}中的目标具有任务效用单位，每项约束则保留原资源单位，因此不存在把小时、金额和效用直接相减的问题。若业务确实要比较净效用，必须预先给出换算系数$\beta_j$，使$\beta_j C_j(a)$与$U$同量纲，再从期望效用增量中减去$\sum_j\beta_j C_j(a)$。期望收益与换算系数都需要根据历史审计、来源覆盖、依赖模型和业务取舍估计；该式提供的是比较框架，不是无需数据就能算出的真实最优策略。

沿用表~\ref{tab:knowledge-fusion-source-records}，对$x_5$再次运行$R_1$不会增加新信息；补充独立人工判断可能区分共享规则与现有人工异议；开发一条新规则有一次性成本，却可批量覆盖同类样本。完整预算还包括规则维护、模型调用延迟、结果审计和失败重试。低单次成本不等于低有效信息成本，批量复用价值也不能替代独立质量验证。

停止条件应与用途对应：预期收益不再超过成本时停止追加，证据达到接收标准时进入使用，关键样本无法自动解决时转交人工，预算不足时保留弃权。分批试用新来源可以先估计互补性，再决定是否扩大投入。跨来源资源分配解决的是当前融合缺口；人工标注内部如何安排新增与重复判断，则需要结合相应的数据获取策略。

\section{本章小结}
\label{sec:knowledge-fusion-summary}

多份知识形成综合结果，依赖的不是一个更复杂的聚合器，而是一条条件清楚的处理链。来源可靠性回答判断是否可信，混淆矩阵保留错误方向，来源相关性防止共享证据被重复计数；对象对齐先确定记录是否指向同一对象，语义对齐确认它们是否回答同一命题，粒度与形式对齐再说明哪些转换成立以及会丢失什么。只有这些条件满足后，投票、加权和概率推断的结果才具有明确含义。

对共同类别答案，等权投票提供可复核基线，加权方法显式区分来源贡献，Dawid--Skene模型与标签模型则把潜在标签、来源误差、弃权和依赖写入生成过程。无参考估计依赖重叠、条件独立或正确建模的依赖结构、方向锚点与秩条件；EM收敛、参数可识别和答案正确是三个不同要求。多个局部答案还受层级、互斥或序列结构约束，必须在局部证据与整体可行性之间联合推断。

融合结果可以硬化为单一标签，也可以保留概率、候选集合或弃权；含义不同的关系、局部目标和表示目标则应通过多个损失或约束联合使用，不能伪装成同一种标签。回看表~\ref{tab:knowledge-fusion-source-records}，共享模板使票数夸大，对象错配会制造虚假冲突，证据不足则应留下未决状态。保留来源、对齐、参数与使用规则的版本，才能继续检查哪些结果可以采用、哪些需要追加信息；这些综合结果是否真的可靠，还要由独立质量评估回答。
